\section{Discussion}

\subsection{Key Findings Interpretation}

The 5$\times$ drift slope difference between CAB and MOHAWK suggests a fundamental mechanistic distinction in what each objective transfers during distillation.

\textbf{We hypothesize that token-level objectives are position-agnostic}: Aligning Q/K projections to B/C parameters would create supervision that doesn't depend on specific sequence positions. If each token's representation is aligned independently, this could produce invariance to sequence length by construction.

\textbf{Matrix-level objectives encode positional relationships}: Attention maps capture which positions attend to which. These position-position patterns are intrinsically length-dependent---a pattern learned at 512 tokens does not describe the same attention structure at 2048 tokens.

The bounded drift ratio (1.34) for CAB versus higher ratio (2.02) for MOHAWK provides evidence consistent with this hypothesized mechanistic difference.

\subsection{Honest Limitations}

We acknowledge several limitations that bound the claims we can make:

\textbf{PoC training only}: Our experiments use 1M tokens (PoC mode) rather than the 1.5B tokens per condition required for full distillation convergence. While drift measurements are valid for comparing representation stability, downstream task performance (F1 on LongBench) cannot be verified.

\textbf{Single teacher model}: All experiments use Phi-1.5 as teacher. Generalization to other Transformers (Llama, Mistral, GPT) is not tested.

\textbf{Simulated student approximation}: Drift analysis uses projection-matched Mamba as student proxy rather than fully-trained phi-mamba checkpoints. The 5$\times$ slope ratio is directionally robust but exact magnitudes may vary with real checkpoints.

\textbf{Context limit discovered}: Phi-1.5's hard 2048 token limit prevented testing at 16K--32K as originally planned. Our results apply to the $\leq$2048 range and suggest (but don't prove) behavior at extrapolated lengths.

\subsection{Broader Impact}

This work provides a principled criterion for distillation objective selection: when target deployment involves sequences longer than teacher training length, token-level objectives produce more stable representations.

For practitioners converting Transformers to SSMs for long-context applications, our 5$\times$ stability gap suggests token-level alignment (CAB-style) may be preferred despite comparable performance at training lengths.
